\documentclass[11pt]{article}
\usepackage[margin=1in]{geometry}
\usepackage{booktabs,amsmath,amssymb}
\usepackage[T1]{fontenc}
\usepackage[hidelinks]{hyperref}

\title{PO-risk Shapley:\\
interpreting distribution shift through the pseudo-outcome model}
\author{SAPerm --- online PO-risk}
\date{}

\begin{document}
\maketitle

\begin{abstract}
\noindent
Predictive SHAP attributes a fitted outcome model $f(X)$ to named
covariates. We apply the same attribution to the doubly robust
pseudo-outcome learner $\tau(X)$ that underlies online PO-risk.
The OOD score remains $\mathbb{E}[\tau(X)^2]$; Shapley values of
$\tau$ name which coordinates drive that score.
On $B{=}20$ linear streams, a pure concept shift lights PO-risk Shapley
(driver hit rate $1.00$) while RF-domain / MMD stay near chance; a pure
covariate shift does the opposite (RF / MMD hit rate $1.00$ on the
shifted coordinate). Same engine as predictive SHAP; only the supervised
target is the DR residual.
\end{abstract}

\section{From predictive SHAP to PO-risk SHAP}

In the predictive setting one fits $f$ for $\mathbb{E}[Y\mid X]$ and
reports Shapley values $\phi_j(x)$ of $f$, with global importance
$\overline{\phi}_j=\mathbb{E}_x|\phi_j(x)|$. That construction needs
only a fitted map $x\mapsto\widehat{g}(x)$.

Online PO-risk uses a different map. With batch indicator
$T\in\{0,1\}$ (reference vs.\ queried window), fit nuisances
$\mu_t(x)=\widehat{\mathbb{E}}[Y\mid X=x,T=t]$ and
$e(x)=\widehat{\mathbb{P}}(T=1\mid X=x)$, form
\begin{equation}
\label{eq:phi}
\phi
=
(Y-\mu_T(X))\,\frac{T-e(X)}{e(X)\,(1-e(X))}
+(\mu_1(X)-\mu_0(X)),
\end{equation}
and regress $\tau(X)\approx\phi$. The scalar OOD score is
$\mathbb{E}[\tau(X)^2]$; PO-risk is $\mathbb{E}[(\phi-\tau(X))^2]$.

\paragraph{Translation rule.}
Replace $f$ by $\tau$. Keep LinearSHAP (ridge) or TreeSHAP (RF).
Global OOD importance is $\overline{\phi}_j=\mathbb{E}|\mathrm{SHAP}_j(\tau;x)|$.
Residual-permutation VIMP (permute coordinate $j$, measure the rise in
PO-risk) is a second view of the same object.

\section{Trichotomy}

\begin{center}
\begin{tabular}{lll}
\toprule
Method & Target & Responds to \\
\midrule
MMD-LOCO / RF-domain & $P(X)$ gap & covariate movement \\
PO-risk Shapley / VIMP & $\tau(X)\approx\phi$ & concept / label-channel residual \\
Online $\tau^2$ score & scalar alarm & either, after nuisances \\
\bottomrule
\end{tabular}
\end{center}

\section{Design}

Stream $Y=X_{\mathrm{sig}}\beta+\varepsilon$, $p{=}8$, $n{=}280$,
change at $t{=}140$, $B{=}20$ seeds.
\texttt{concept}: drop $\beta$ on $\{2,3\}$ after the change ($P(X)$ fixed).
\texttt{covariate}: add a constant to $x_0$ ($Y\mid X$ fixed).
Each seed: one pre window (two halves of the burn-in) and one post
window (burn-in vs.\ post-shift block of size $50$).
$\tau\in\{\mathrm{RF},\mathrm{ridge}\}$, $N_{\mathrm{perm}}{=}15$.
Sliding windows: $w{=}60$, step $20$, RF $\tau$.
Code: \texttt{python -m experiments.online\_po\_risk.run\_po\_shapley --seeds 20}.

\section{Results}

\subsection{Window-level recovery}

\begin{table}[ht]
\centering
\caption{$B{=}20$. Driver hit rate $=$ fraction of seeds where
$\arg\max_j$ lands in the designed driver set
($\{2,3\}$ for concept, $\{0\}$ for covariate).
Mean driver vs.\ noise mass for each attribution.}
\label{tab:window}
\scriptsize
\begin{tabular}{ll l ccc cc cc}
\toprule
Scenario & Phase & $\tau$
& \multicolumn{3}{c}{top hit rate}
& \multicolumn{2}{c}{Shapley}
& \multicolumn{2}{c}{RF-domain} \\
\cmidrule(lr){4-6}\cmidrule(lr){7-8}\cmidrule(lr){9-10}
& & & Shap & RF & MMD & driver & noise & driver & noise \\
\midrule
concept & pre  & RF    & .20 & .20 & .05 & .10 & .12 & .13 & .13 \\
concept & post & RF    & \textbf{1.00} & .35 & .05 & \textbf{.65} & .12 & .13 & .12 \\
concept & pre  & ridge & .20 & .20 & .05 & .14 & .16 & .13 & .13 \\
concept & post & ridge & \textbf{1.00} & .35 & .05 & \textbf{.91} & .15 & .13 & .12 \\
\midrule
covariate & pre  & RF    & .10 & .10 & .45 & .12 & .12 & .12 & .13 \\
covariate & post & RF    & .20 & \textbf{1.00} & \textbf{1.00} & .14 & .12 & \textbf{.37} & .09 \\
covariate & pre  & ridge & .20 & .10 & .45 & .20 & .17 & .12 & .13 \\
covariate & post & ridge & .20 & \textbf{1.00} & \textbf{1.00} & .24 & .19 & \textbf{.37} & .09 \\
\bottomrule
\end{tabular}
\end{table}

Post-concept $\tau^2$ rises from $0.41$ to $2.63$ (RF) / $4.29$ (ridge);
post-covariate $\tau^2$ stays near $0.5$. The alarm and the Shapley
localization move together on concept drift; on covariate drift the
alarm is quiet and $P(X)$-tools take the localization.

\subsection{Feature-level means (post, RF $\tau$)}

\begin{table}[ht]
\centering
\caption{Mean attribution across $B{=}20$ post windows (RF $\tau$).
Concept: Shapley / VIMP peak on $\{x_2,x_3\}$.
Covariate: RF-domain / MMD peak on $x_0$.}
\label{tab:feat}
\scriptsize
\begin{tabular}{l rrrr c | l rrrr c}
\toprule
\multicolumn{6}{c|}{concept} & \multicolumn{6}{c}{covariate} \\
\cmidrule(lr){1-6}\cmidrule(lr){7-12}
Feat & Shap & VIMP & RF & MMD & drv
& Feat & Shap & VIMP & RF & MMD & drv \\
\midrule
$x_0$ & .11 & .27 & .12 & .00 &
& $x_0$ & .14 & .33 & \textbf{.37} & \textbf{.05} & yes \\
$x_1$ & .16 & .37 & .11 & .00 &
& $x_1$ & .10 & .28 & .10 & $-$.01 & \\
$x_2$ & \textbf{.62} & \textbf{1.93} & .13 & .00 & yes
& $x_2$ & .11 & .27 & .09 & $-$.01 & \\
$x_3$ & \textbf{.67} & \textbf{2.27} & .13 & .00 & yes
& $x_3$ & .12 & .28 & .09 & $-$.01 & \\
$x_4$ & .09 & .26 & .13 & .00 &
& $x_4$ & .10 & .26 & .09 & $-$.01 & \\
$x_5$ & .13 & .32 & .13 & .00 &
& $x_5$ & .12 & .33 & .09 & $-$.01 & \\
$x_6$ & .09 & .26 & .12 & .00 &
& $x_6$ & .13 & .32 & .08 & $-$.01 & \\
$x_7$ & .11 & .27 & .13 & .00 &
& $x_7$ & .15 & .35 & .09 & $-$.01 & \\
\bottomrule
\end{tabular}
\end{table}

\subsection{Sliding windows}

\begin{table}[ht]
\centering
\caption{Online Shapley of $\tau$ (RF), $w{=}60$, step $20$, $B{=}20$.
After a concept change, mean driver Shapley rises and the top-1 driver
rate nearly doubles; covariate windows stay flat on Shapley while the
fixed-window RF/MMD analysis above already localizes $x_0$.}
\label{tab:slide}
\begin{tabular}{ll cccc}
\toprule
Scenario & Phase & top-driver rate & driver Shap & noise Shap & $\tau^2$ \\
\midrule
concept & pre  & .26 & .11 & .12 & .44 \\
concept & post & \textbf{.44} & \textbf{.20} & .12 & \textbf{.74} \\
covariate & pre  & .12 & .11 & .12 & .44 \\
covariate & post & .10 & .12 & .11 & .47 \\
\bottomrule
\end{tabular}
\end{table}

\section{Operational reading}

\begin{enumerate}
\item Alarm on $\mathbb{E}[\tau^2]$.
\item On a rejecting window, read $\overline{\phi}_j$ of $\tau$ for
  concept drivers.
\item If RF-domain / MMD also spike, the window mixes covariate mass;
  do not claim a pure concept mechanism.
\item The ranking is the same deliverable predictive SHAP already
  gave for $f(X)$---only the label is the DR residual.
\end{enumerate}

\section*{Artifacts}
\texttt{po\_shapley\_window.csv},
\texttt{po\_shapley\_window\_summary.csv},
\texttt{po\_shapley\_feature\_means.csv},
\texttt{online\_po\_shapley.csv},
\texttt{online\_po\_shapley\_summary.csv}
under \texttt{experiments/online\_po\_risk/results/}.

\section*{Takeaway}
PO-risk $+$ Shapley is predictive SHAP with the label replaced by the
DR batch residual. Detection stays on $\tau^2$; interpretation stays
on SHAP of $\tau$. Concept drift is exactly where that pair works and
where $P(X)$-only tools are silent; covariate drift is the complementary
case.

\end{document}
